\chapter{Метод стрельбы}
        
Для нахождения устойчивой переходной траектории используется метод стрельбы. Его идея состоит в следующем: при фиксированном начальном значении капитала $K_0$ подбирается такое начальное значение потребления $C_0$, при котором траектория системы на конечном горизонте максимально приближается к стационарному состоянию.

На каждом шаге алгоритма выполняются следующие действия:
\begin{enumerate}[label=\arabic*)]
    \item по заданным $K_t$ и $C_t$ вычисляется $L_t$;
    \item определяется выпуск $Y_t$ и инвестиции $I_t=Y_t-C_t$;
    \item по уравнению накопления вычисляется капитал $K_{t+1}$;
    \item численно решается нелинейное уравнение относительно $C_{t+1}$:
    \begin{equation}
    C_{t+1}=\beta\bigl(1+R(K_{t+1},C_{t+1})-\delta\bigr)C_t;
    \end{equation}
    \item процесс повторяется до достижения заданного горизонта моделирования.
\end{enumerate}

Подбор начального потребления осуществляется по минимуму функционала, измеряющего отклонение терминального состояния от стационарной точки.
    
\section{Построение траектории при фиксированном \texorpdfstring{$C_0$}{C0}}

    Пусть заданы начальные значения $K_0$ и $C_0$. Для каждого периода $t=0,1,\dots,T-1$ траектория строится итерационно. Сначала по внутрипериодному условию вычисляется труд:
    \[
    L_t=\left(\frac{(1-\alpha)A}{\phi}\frac{K_t^\alpha}{C_t}\right)^{\frac{1}{\psi+\alpha}}.
    \]
    Затем определяются выпуск и связанные с ним величины:
    \(Y_t=A K_t^\alpha L_t^{1-\alpha}\),
    \(R_t=\alpha\frac{Y_t}{K_t}\),
    \(I_t=Y_t-C_t\).
    После этого капитал следующего периода находится из закона накопления:
    \(K_{t+1}=Y_t-C_t+(1-\delta)K_t\).

    Значение $C_{t+1}$ определяется из уравнения Эйлера и находится как корень одномерного нелинейного уравнения
    \[
    f(c)=c-\beta\bigl(1+R(K_{t+1},c)-\delta\bigr)C_t=0,
    \]
    где \(c=C_{t+1}\).
    При этом функция доходности на шаге \(t+1\) задается через
    \(Y_{t+1}=A K_{t+1}^\alpha L_{t+1}^{1-\alpha}\),
    \(L_{t+1}=\left(\frac{(1-\alpha)A}{\phi}\frac{K_{t+1}^\alpha}{c}\right)^{\frac{1}{\psi+\alpha}}\),
    \(R_{t+1}=\alpha\frac{Y_{t+1}}{K_{t+1}}\).
    После нахождения \(C_{t+1}\) выполняется переход к периоду \(t+1\), и процедура повторяется до горизонта \(T\).
\section{Метод Ньютона для нахождения \texorpdfstring{$C_{t+1}$}{C(t+1)}}

    После вычисления \(K_{t+1}\) значение \(C_{t+1}\) определяется из неявного уравнения Эйлера:
    \[
    f(c)=c-\beta\bigl(1+R(K_{t+1},c)-\delta\bigr)C_t=0,\qquad c=C_{t+1}.
    \]
    Для фиксированного \(K=K_{t+1}\) используем обозначения
    \(L(K,c)=\left(\frac{(1-\alpha)A}{\phi}\frac{K^\alpha}{c}\right)^{\frac{1}{\psi+\alpha}}\),
    \(Y(K,c)=A K^\alpha L(K,c)^{1-\alpha}\),
    \(R(K,c)=\alpha\frac{Y(K,c)}{K}\).

    Производная \(f'(c)\) вычисляется аналитически. Из
    \(\frac{\partial L}{\partial c}=-\frac{1}{\psi+\alpha}\frac{L}{c}\),
    \(\frac{\partial Y}{\partial c}=-\frac{1-\alpha}{\psi+\alpha}\frac{Y}{c}\)
    получаем
    \(\frac{\partial R}{\partial c}=-\frac{1-\alpha}{\psi+\alpha}\frac{R}{c}\), а значит
    \[
    f'(c)=1-\beta C_t\frac{\partial R(K,c)}{\partial c}
    =1+\beta C_t\frac{1-\alpha}{\psi+\alpha}\frac{R(K,c)}{c}.
    \]

    Тогда итерации Ньютона записываются в виде
    \[
    c^{(n+1)}=
    c^{(n)}-\frac{f(c^{(n)})}{f'(c^{(n)})}
    =
    c^{(n)}-
    \frac{c^{(n)}-\beta\bigl(1+R(K_{t+1},c^{(n)})-\delta\bigr)C_t}
    {1+\beta C_t \frac{1-\alpha}{\psi+\alpha}\frac{R(K_{t+1},c^{(n)})}{c^{(n)}}}.
    \]
    После достижения критерия сходимости принимается \(C_{t+1}=c^{(n)}\).

\section{Псевдокод метода стрельбы}
    \begin{enumerate}[label=\arabic*)]
        \item Задать параметры модели $(\alpha,\beta,\delta,\phi,\psi,A)$.
        \item Вычислить стационарное состояние $(K^*,C^*,L^*)$.
        \item Зафиксировать $K_0$.
        \item Задать начальный интервал для потребления: $C_0\in[C_{\min},C_{\max}]$.
        \item Повторять внешнюю бисекцию по $C_0$, пока $|C_{\max}-C_{\min}|\ge\varepsilon_{\text{outer}}$.
        \item На каждой итерации положить $C_0=\frac{C_{\min}+C_{\max}}{2}$.
        \item Построить траекторию на горизонте $t=0,1,\dots,T-1$.
        \item Для каждого $t$ вычислить \(L_t=\left(\frac{(1-\alpha)A}{\phi}\frac{K_t^\alpha}{C_t}\right)^{\frac{1}{\psi+\alpha}}\).
        \item Для каждого $t$ вычислить \(Y_t=A K_t^\alpha L_t^{1-\alpha}\).
        \item Для каждого $t$ вычислить \(K_{t+1}=Y_t-C_t+(1-\delta)K_t\).
        \item Для каждого $t$ найти $C_{t+1}$ методом Ньютона из \(f(c)=c-\beta\bigl(1+R(K_{t+1},c)-\delta\bigr)C_t=0\).
        \item По завершении траектории вычислить терминальное отклонение $g(C_0)=K_T-K^*$.
        \item Если $g(C_0)>0$, положить $C_{\min}=C_0$.
        \item Если $g(C_0)\le0$, положить $C_{\max}=C_0$.
        \item После сходимости внешней бисекции принять найденное $C_0^*$.
        \item Построить итоговую saddle-path траекторию $(K_t,C_t,L_t)$.
    \end{enumerate}
